% Computer Science A --- Object-oriented Java (inheritance, polymorphism, abstract classes, interfaces).
% \input from \texttt{main.tex} before \texttt{cs-a/dsa.tex}.
% Design-pattern FRQs (Factory, Singleton) migrated from \texttt{cs-a/05-practice.tex} appear at the end of this file.

\runningheader{Computer Science A}{Object-oriented Java}{Page \thepage\ of \numpages}

\uplevel{\par\textbf{Classes, objects, and memory}\par}

\question[4]
What are classes used for?
\begin{solutionorbox}[1.2in]
  Blueprints for objects: they group fields (state) and methods (behavior), define types, encapsulate implementation, and support reuse and modeling of real-world concepts.
\end{solutionorbox}

\question[4]
Describe what an object is.
\begin{solutionorbox}[1.15in]
  A runtime instance of a class: a specific chunk of state in memory with a reference identity, whose behavior is defined by the class's methods.
\end{solutionorbox}

\question[4]
What is the role of garbage collection in Java?
\begin{solutionorbox}[1.2in]
  The JVM reclaims memory occupied by objects that are no longer reachable, so programmers rarely \texttt{free} manually; it reduces leaks and dangling pointers at the cost of pauses and tuning complexity.
\end{solutionorbox}

\question[4]
Describe the difference between the stack and the heap.
\begin{solutionorbox}[1.45in]
  The stack holds frames for active method calls (locals, parameters, partial execution state), typically LIFO and size-known; the heap stores objects and grows dynamically; references on the stack often point to heap objects.
\end{solutionorbox}

\question[4]
What is a constructor and how is it different from a method?
\begin{solutionorbox}[1.35in]
  A constructor initializes a new instance; its name matches the class, it has no return type, and it runs when \texttt{new} is used. Ordinary methods can have any name, a return type (or \texttt{void}), and are invoked on an instance (or statically).
\end{solutionorbox}

\newpage

\uplevel{\par\textbf{Inheritance and polymorphism}\par}

\question[4]
Describe inheritance and how would you use it in a project?
\begin{solutionorbox}[1.55in]
  A subclass extends a superclass, inheriting fields/methods and adding or overriding behavior. Use it to model ``is-a'' relationships, share common implementation, and enable polymorphic references to superclass types.
\end{solutionorbox}

\question[4]
Describe polymorphism and how would you use it in a project?
\begin{solutionorbox}[1.5in]
  Same interface or superclass type, different runtime behavior: overridden methods dispatch on the actual object type. Use it to write generic code (e.g.\ process a \texttt{List} of \texttt{Shape} without branching on concrete classes).
\end{solutionorbox}

\question[4]
How is method overriding different from method overloading?
\begin{solutionorbox}[1.25in]
  Overriding replaces a superclass \emph{instance} method implementation with the same signature in a subclass (runtime dispatch). Overloading defines multiple methods with the \emph{same name} but different parameter lists in one class (compile-time resolution).
\end{solutionorbox}

\question[4]
What methods are commonly overridden from the \texttt{Object} class and why?
\begin{solutionorbox}[1.45in]
  \texttt{equals}/\texttt{hashCode} for value equality and hash-based collections; \texttt{toString} for debugging and logging; sometimes \texttt{clone} or \texttt{finalize} (rare today). Contracts must stay consistent (e.g.\ equal objects \(\Rightarrow\) same hash).
\end{solutionorbox}

\newpage

\uplevel{\par\textbf{Designing class hierarchies}\par}

\question[4]
What state and behavior would you define in a parent class but not a subclass?
\begin{solutionorbox}[1.45in]
  Shared data and logic common to all specializations (e.g.\ name/ID for all employees, generic validation), while keeping the type general enough that subclasses add specifics rather than duplicating the common core.
\end{solutionorbox}

\question[4]
What state and behavior would you define in a subclass but not in the parent class?
\begin{solutionorbox}[1.45in]
  Details only for that subtype (hourly wage vs.\ annual salary, etc.), so the base class is not cluttered with fields irrelevant to every sibling.
\end{solutionorbox}

\question[4]
How does upcasting an object work?
\begin{solutionorbox}[1.15in]
  Assign a subclass reference to a superclass type (implicit); the object stays the same; only the compile-time type narrows what members are visible without a cast.
\end{solutionorbox}

\question[4]
How does downcasting an object work?
\begin{solutionorbox}[1.2in]
  Use an explicit cast from superclass to subclass when you know the runtime type; it can throw \texttt{ClassCastException} if wrong. Prefer \texttt{instanceof} (pattern matching) before casting when uncertain.
\end{solutionorbox}

\newpage

\uplevel{\par\textbf{Abstraction, encapsulation, and access}\par}

\question[4]
Describe abstraction and how would you use it in a project?
\begin{solutionorbox}[1.45in]
  Hiding complexity behind simple interfaces (abstract types, APIs, layers). Use abstract classes/interfaces to express ``what'' without ``how,'' so implementations can change without affecting clients.
\end{solutionorbox}

\question[4]
Describe encapsulation and how would you use it in a project?
\begin{solutionorbox}[1.45in]
  Hide internal state behind methods (private fields, public accessors/mutators) so invariants and validation stay centralized and dependencies on representation do not leak.
\end{solutionorbox}

\question[4]
What are the four levels of access we can give to class members? How are they different from one another?
\begin{solutionorbox}[1.55in]
  \texttt{private} (enclosing class), package/default, \texttt{protected} (package + subclasses), \texttt{public}. Narrower access hides implementation details.
\end{solutionorbox}

\question[4]
What is the purpose of using getter and setter methods?
\begin{solutionorbox}[1.25in]
  Control reads/writes to fields (validation, lazy init, logging), preserve encapsulation, and allow evolving internal storage without breaking callers.
\end{solutionorbox}

\newpage

\uplevel{\par\textbf{Interfaces and abstract types}\par}

\question[4]
Why would you use an interface over an abstract class?
\begin{solutionorbox}[1.35in]
  Java allows multiple interface implementation but single inheritance; interfaces define capability contracts without state (pre--Java~8) or with default/static methods; good for cross-cutting roles and test doubles.
\end{solutionorbox}

\newpage

\begingradingrange{csoo}

\uplevel{\par\textbf{Object-Oriented Programming and Java}\par}

\question[4]
True or false: Interface implementation is a form of inheritance.
\begin{choices}
  \CorrectChoice FALSE
  \choice TRUE
\end{choices}
\begin{solution}
  In Java, \texttt{extends} (class inheritance) and \texttt{implements} (interface contract) are different: implementing an interface is fulfilling a contract, not class inheritance in the \texttt{extends} sense.

  \begin{center}
  \small
  \setlength{\tabcolsep}{5pt}
  \begin{tabular}{@{}p{0.18\linewidth}p{0.38\linewidth}p{0.38\linewidth}@{}}
  \hline
  \textbf{Aspect} & \textbf{Inheritance (\texttt{extends})} & \textbf{Implementation (\texttt{implements})} \\
  \hline
  Relationship & \texttt{is-a} & \texttt{acts-as} (fulfills a contract) \\
  Code reuse & High: you inherit method bodies & Low: you must write your own method bodies \\
  Limit & Only one superclass & Many interfaces allowed \\
  Focus & Shared identity and internal logic & External behavior (what others can call) \\
  \hline
  \end{tabular}
  \end{center}
\end{solution}

\question[4]
The relation between a superclass which gets extended by a single subclass is called \blank[4em]\ inheritance.
\begin{choices}
  \choice Multilevel inheritance
  \CorrectChoice Single inheritance
  \choice Hybrid inheritance
  \choice Multiple inheritance
\end{choices}
\begin{solution}
  One subclass extending one direct superclass is \textbf{single inheritance}. Java classes support single inheritance of implementation; multiple interfaces are separate.
\end{solution}

\question[4]
The keyword \texttt{super} calls the constructor of the \blank[6em].
\begin{choices}
  \CorrectChoice Parent class
  \choice Subclass
  \choice Current class
  \choice Specified class
\end{choices}
\begin{solution}
  \texttt{super(\ldots)} in a constructor chains to the superclass constructor; \texttt{super.method()} calls a superclass instance method.
\end{solution}

\question[4]
What is the purpose of the \texttt{super} keyword in Java?
\begin{choices}
  \choice To represent an instance of the superclass
  \choice To call a static method
  \CorrectChoice To invoke the superclass's constructor
  \choice To access the current object's properties
\end{choices}
\begin{solution}
  \texttt{super(\ldots)} chains to the superclass constructor; \texttt{super} can also qualify members hidden by the subclass, but among these options the constructor call is the intended answer.
\end{solution}

\newpage

\question[4]
What is the output of the following code?
\begin{lstlisting}[style=java]
import java.util.*;
public class Olympics {
    public static void main(String[] args) {
        India i = new India();
        i.hostOlympics();
    }
    public void hostOlympics() {
        System.out.println("Summer Olympics and Winter Olympics");
    }
}
class India extends Olympics {
    public void hostOlympics() {
        System.out.println("We will host the summer olympics in 2036");
        Olympics o = new Olympics();
        o.hostOlympics();
    }
}
class Japan extends Olympics {
    public void hostOlympics() {
        System.out.println("We hosted Olympics on 2021");
    }
}
\end{lstlisting}
\begin{choices}
  \choice \texttt{Summer Olympics and Winter Olympics} (then \texttt{We will host the summer olympics in 2036})
  \choice \texttt{Summer Olympics and Winter Olympics}
  \CorrectChoice \texttt{We will host the summer olympics in 2036} then \texttt{Summer Olympics and Winter Olympics}
  \choice \texttt{We will host the summer olympics in 2036}
\end{choices}
\begin{solution}
  \texttt{India.hostOlympics} runs the override, prints the India line, then calls \texttt{new Olympics().hostOlympics()}, which uses the superclass method and prints the second line.
\end{solution}

\question[4]
Which of the following is true when an abstract class implements an interface?
\begin{choices}
  \choice Abstract class must override all the methods of an interface
  \choice None of the above
  \CorrectChoice The abstract class does not need to override all the methods of the interface
  \choice Abstract class must override all the abstract methods of an interface
\end{choices}
\begin{solution}
  An abstract class may leave some interface methods abstract for concrete subclasses to implement.
\end{solution}

\question[4]
What is the output of the following Java code?
\begin{lstlisting}[style=java]
class Animal {
    public void barkingDog() {
        System.out.println("Wolf Wolf!");
    }
    public void barkingDog(String a) {
        System.out.println(a);
    }
    public static void main(String[] args) {
        Animal a = new Animal();
        a.barkingDog("Barking!");
    }
}
\end{lstlisting}
\begin{choices}
  \choice Compile-time error
  \choice None of the above
  \choice Run-time error
  \CorrectChoice \texttt{Barking!}
\end{choices}
\begin{solution}
  The call matches the overloaded \texttt{barkingDog(String)} method, which prints its argument.
\end{solution}

\newpage

\question[4]
What members are allowed in an abstract class?
\begin{choices}
  \choice Constructor
  \choice Non-abstract method
  \choice Abstract method
  \CorrectChoice All of the above
\end{choices}
\begin{solution}
  Abstract classes may contain constructors, concrete methods, fields, and \texttt{abstract} methods; they cannot be instantiated directly.
\end{solution}

\question[4]
Garbage collection: when is the object created on the first line eligible for garbage collection?
\begin{lstlisting}[style=java]
public static void main(String[] args) {
    Object obj1 = new Object(); // 1st line
    Object obj2 = obj1; // 2nd line
    Object obj3 = new Object(); // 3rd line
    obj1 = obj3; // 4th line
    obj2 = obj3; // 5th line
    obj1 = new Object(); // 6th line
}
\end{lstlisting}
\begin{choices}
  \CorrectChoice After the 5th line
  \choice After the 6th line
  \choice After the 4th line
  \choice After the 3rd line
\end{choices}
\begin{solution}
  After line~5, both \texttt{obj1} and \texttt{obj2} refer to \texttt{obj3}; the original object from line~1 has no references and may be collected.
\end{solution}

\question[4]
What is runtime polymorphism?
\begin{choices}
  \CorrectChoice Runtime polymorphism is a process in which a call to an \textbf{overridden} method is resolved at runtime rather than at compile-time
  \choice None of the above
  \choice Both of the above
  \choice Runtime polymorphism is a process in which a call to an \textbf{overloaded} method is resolved at runtime rather than at compile-time
\end{choices}
\begin{solution}
  \textbf{Overloading} is compile-time (static) polymorphism; \textbf{overriding} uses dynamic dispatch at runtime.
\end{solution}

\question[4]
Which keyword is used for class inheritance in Java?
\begin{choices}
  \CorrectChoice \texttt{extends}
  \choice \texttt{:}
  \choice \texttt{is}
  \choice \texttt{implements}
\end{choices}
\begin{solution}
  \texttt{extends} declares a subclass; \texttt{implements} lists interfaces.
\end{solution}

\newpage

\question[4]
You are working on an application that throws an exception when it runs. What should be your \emph{first} step to debug the issue?
\begin{choices}
  \choice Change the code at random
  \choice Comment out parts of the code without reading the error
  \choice Ask your instructor immediately without gathering details
  \CorrectChoice Read the stack trace (and message) from the exception
\end{choices}
\begin{solution}
  The stack trace shows which exception was thrown and the call chain (classes, methods, and line numbers) so you can locate the failure. Then you can form a hypothesis, reproduce, and fix; blind edits or skipping the error output waste time.
\end{solution}

\question[4]
What will be the output when the following program is run? (Assume line numbers in the messages below are as given.)
\begin{lstlisting}[style=java]
package exceptions;
public class TestClass {
    public static void main(String[] args) {
        try {
            hello();
        } catch (MyException me) {
            System.out.println(me);    // short message
            // me.printStackTrace();   // full stack trace
        }
    }
    static void hello() throws MyException {
        int[] dear = new int[7];
        dear[0] = 747;
        foo();
    }
    static void foo() throws MyException {
        throw new MyException("Exception from foo");
    }
}
class MyException extends Exception {
    public MyException(String msg) {
        super(msg);
    }
}
\end{lstlisting}
\begin{choices}
  \choice Full stack trace from \texttt{printStackTrace}-style output
  \CorrectChoice \texttt{exceptions.MyException: Exception from foo} (no stack trace, because only \texttt{println(me)} runs)
  \choice \texttt{ArrayIndexOutOfBoundsException} on the array line
  \choice \texttt{Error in thread "main"} with \texttt{ArrayIndexOutOfBoundsException}
\end{choices}
\begin{solution}
  \texttt{System.out.println(me)} prints the exception's short string (\texttt{toString()}); the stack trace is commented out. \texttt{foo} throws \texttt{MyException} as shown.
\end{solution}

\question[4]
Constructors have no return type.
\begin{choices}
  \CorrectChoice TRUE
  \choice FALSE
\end{choices}
\begin{solution}
  Constructor names match the class name and they do not declare a return type; \texttt{void} would make them ordinary methods, not constructors.
\end{solution}

\question[4]
A constructor is \blank[10em].
\begin{choices}
  \choice A normal method
  \CorrectChoice A method called when a new object is created using the \texttt{new} keyword
  \choice A field which is of the same type as the class it is defined within
  \choice A way of implementing encapsulation by providing a method for mutating values
\end{choices}
\begin{solution}
  Constructors initialize new instances; they often establish invariants before the object is used.
\end{solution}

\newpage

\question[4]
What will be the output when the following program is run? (Assume there is no error in the line numbers given in the options.)
\begin{lstlisting}[style=java]
package exceptions;
public class TestClass {
    public static void main(String[] args) {
        try {
            doTest();
        } catch (MyException me) {
            System.out.println(me);
        }
    }
    static void doTest() throws MyException {
        int[] array = new int[10];
        array[10] = 1000;
        doAnotherTest();
    }
    static void doAnotherTest() throws MyException {
        throw new MyException("Exception from doAnotherTest");
    }
}
class MyException extends Exception {
    public MyException(String msg) {
        super(msg);
    }
}
\end{lstlisting}
\begin{choices}
  \choice Caught \texttt{MyException} from \texttt{doAnotherTest}
  \CorrectChoice Uncaught \texttt{ArrayIndexOutOfBoundsException} at \texttt{array[10]} (the \texttt{try} never completes successfully)
  \choice \texttt{Error in thread "main"} with \texttt{ArrayIndexOutOfBoundsException}
  \choice \texttt{exceptions.MyException: Exception from doAnotherTest} only
\end{choices}
\begin{solution}
  Assigning to \texttt{array[10]} on an array of length~10 throws \texttt{ArrayIndexOutOfBoundsException}, a \texttt{RuntimeException} subclass. It is not caught by \texttt{catch (MyException)}, so it propagates and the default thread handler prints the stack trace.
\end{solution}

\question[4]
Which is an excellent example of polymorphism? (Select all that apply.)
\begin{checkboxes}
  \CorrectChoice Overloading methods (compile-time polymorphism)
  \choice Extending a parent class
  \choice Transpiling TypeScript into JavaScript
  \choice Adding elements to a collection
\end{checkboxes}
\begin{solution}
  Overloading is the usual ``ad-hoc'' polymorphism example at compile time. The other distractors are not polymorphism in the OOP sense.
\end{solution}

\question[4]
A department uses one 3D printer to print both lab equipment and a toy action figure from different digital models. This illustrates which idea?
\begin{choices}
  \choice Abstraction
  \choice Inheritance
  \choice Encapsulation
  \CorrectChoice Polymorphism
\end{choices}
\begin{solution}
  One \texttt{print()} operation behaves differently depending on the model (the ``message''): same interface, different outcomes---an analogy for polymorphism.

  \begin{itemize}
    \item \textbf{Interface}: the \texttt{print()} action.
    \item \textbf{Input}: the blueprint (\textbf{class}/\textbf{object}).
    \item \textbf{Polymorphism}: one operation, many concrete behaviors.
  \end{itemize}
\end{solution}

\newpage

\question[4]
What is the main advantage of inheritance in OOP?
\begin{choices}
  \choice Enforcing encapsulation of data
  \choice Simplifying code by using primitive data types
  \CorrectChoice Reducing code duplication and promoting code reuse
  \choice Enhancing security by hiding the implementation details
\end{choices}
\begin{solution}
  Inheritance lets subclasses reuse and specialize superclass code; encapsulation and primitives are separate concerns.
\end{solution}

\question[4]
John got a recipe for cookies from his mother; John likes chocolate chips, so he altered the original recipe. This is an example of which concept?
\begin{choices}
  \choice Method overloading
  \choice Inheritance
  \choice Abstraction
  \CorrectChoice Method overriding
\end{choices}
\begin{solution}
  Altering inherited behavior while keeping the same ``recipe'' (method signature / role) matches \textbf{overriding} a superclass method in a subclass.
\end{solution}

\question[4]
Which are benefits of polymorphism? (Select all that apply.)
\begin{checkboxes}
  \CorrectChoice It makes the code more reusable
  \CorrectChoice It makes the code more dynamic
  \choice It makes the code more efficient
  \choice It protects the code by preventing extension
\end{checkboxes}
\begin{solution}
  Polymorphism improves structure and flexibility, not raw speed; virtual calls can have small overhead. It \emph{enables} extension (e.g.\ new payment types), unlike \texttt{final}-heavy designs.

  Reusable example: \texttt{processPayment(Payment p)} can accept \texttt{CreditCard}, \texttt{PayPal}, \texttt{Bitcoin}, and future types without duplicating the calling code.

  ``Dynamic'' here means runtime binding: the JVM selects the implementation from the actual object type, so new subclasses can plug in without changing callers.
\end{solution}

\newpage

\uplevel{\par\textbf{Scopes, types, collections, and language rules (practice bank)}\par}

\question[4]
In the code below, \texttt{x} sits in what scope?
\begin{lstlisting}[style=java]
public class Greetings {
    static String w = "welcome";
    String x = "hello";
    void speak() {
        String y = "hi";
        for (int z = 0; z < 4; z++) {
            System.out.println(z);
        }
    }
}
\end{lstlisting}
\begin{choices}
  \choice Final scope
  \choice Block scope
  \choice Global scope
  \CorrectChoice Instance scope
  \choice Static / class scope
  \choice Method scope
\end{choices}
\begin{solution}
  \texttt{x} is an instance field (non-\texttt{static}) of \texttt{Greetings}---\textbf{instance scope}. \texttt{w} is class scope; \texttt{y} and \texttt{z} are narrower (method / loop block).
\end{solution}

\question[4]
\textit{(Select all that apply.)} What is correct about an abstract class?
\begin{checkboxes}
  \choice A concrete class can have multiple abstract classes inherited
  \CorrectChoice In order to inherit an abstract class when using a concrete class, we can use the keyword \texttt{extends}
  \choice An abstract class cannot have any concrete methods predefined
  \CorrectChoice An abstract class does not need to have any predefined methods
\end{checkboxes}
\begin{solution}
  Java allows only one superclass (\texttt{extends}), but that superclass may be abstract. Abstract classes \emph{may} contain concrete methods; they need not declare any abstract methods unless the class itself is abstract and you want to force subclasses to implement something (or the class is not intended to be instantiated).
\end{solution}

\question[4]
\textit{(Select all that apply.)} Assume that \texttt{a}, \texttt{b}, and \texttt{c} refer to instances of wrapper classes. Which of the following statements are correct?
\begin{checkboxes}
  \CorrectChoice \texttt{a.equals(a)} will always return true if their values are the same and if they are instances of the same class
  \choice \texttt{b.equals(c)} may return false even if \texttt{c.equals(b)} returns true
  \choice \texttt{a.equals(b)} returns the same as \texttt{a == b}
  \CorrectChoice \texttt{a.equals(b)} returns false if they refer to instances of different classes
\end{checkboxes}
\begin{solution}
  For wrapper types, \texttt{equals} is value-based for the same runtime type; \texttt{==} compares references. \texttt{equals} should be symmetric for well-behaved types, so the second distractor is not correct in general.
\end{solution}

\question[4]
The following is how we create a class in Java:
\begin{lstlisting}[style=java]
public void buildString(String start, String add) {
    String toBuild = start;
    for (int i = 0; i < 100; i++) {
        toBuild.concat(add);
    }
    System.out.println(toBuild);
}
\end{lstlisting}
\begin{choices}
  \choice TRUE
  \CorrectChoice FALSE
\end{choices}
\begin{solution}
  This is a method, not a class definition. Also, \texttt{String.concat} returns a new \texttt{String}; ignoring the return value means \texttt{toBuild} never changes.
\end{solution}

\question[4]
Which use of the \texttt{final} keyword will lead to compilation errors?
\begin{lstlisting}[style=java]
public final class Example {
    public final void foo() {
        final int x = 4;
        final int y = 3;
        y++;
        System.out.println(x * y);
    }
}
\end{lstlisting}
\begin{choices}
  \choice On the class
  \choice On the method
  \choice On the variable \texttt{x}
  \CorrectChoice On the variable \texttt{y}
\end{choices}
\begin{solution}
  \texttt{y} is \texttt{final} and cannot be incremented. \texttt{final} on the class and method is legal here.
\end{solution}

\question[4]
Which is true about arrays in Java?
\begin{choices}
  \choice Arrays are dynamic and can change size by adding elements
  \CorrectChoice Arrays have a fixed size and cannot grow or shrink
  \choice Arrays are a primitive data type
  \choice Arrays are immutable, so elements inserted into an array cannot be changed in any way
\end{choices}
\begin{solution}
  Array \textbf{length} is fixed at creation; elements are mutable unless the reference is \texttt{final}. Arrays are reference types, not primitives.
\end{solution}

\question[4]
A \texttt{try} block must be followed by a \texttt{catch} block and a \texttt{finally} block.
\begin{choices}
  \choice TRUE
  \CorrectChoice FALSE
\end{choices}
\begin{solution}
  A \texttt{try} must be followed by at least one of \texttt{catch} or \texttt{finally} (or both). \texttt{try} alone is a compile error; \texttt{try} with only \texttt{finally} is legal when you want cleanup without handling exceptions here.
\begin{lstlisting}[style=java]
// try + catch (no finally required)
try {
    int n = Integer.parseInt(s);
} catch (NumberFormatException e) {
    System.out.println("bad number");
}

// try + finally (no catch) --- valid; exceptions propagate after finally runs
try {
    work();
} finally {
    cleanup();
}

// try + catch + finally --- all three allowed together
try {
    readFile();
} catch (IOException e) {
    log(e);
} finally {
    closeQuietly();
}
\end{lstlisting}
\end{solution}

\question[4]
Checked exceptions must \emph{always} be caught in a \texttt{try}-\texttt{catch} block.
\begin{choices}
  \choice TRUE
  \CorrectChoice FALSE
\end{choices}
\begin{solution}
  They must be handled \emph{or} declared with \texttt{throws}; they need not be caught locally.
\end{solution}

\question[4]
Fill in the blank to complete the JUnit test:
\begin{lstlisting}[style=java]
public class Addition {
    public int methodToTest(int x, int y) {
        return x + y;
    }
}
public class AdditionTest {
    @Test
    public void testAddPositiveInts() {
        int n = 2;
        int m = 3;
        _______________________________
    }
}
\end{lstlisting}
\begin{choices}
  \CorrectChoice \texttt{Assert.assertEquals(5, methodToTest(n, m));}
  \choice \texttt{Assert.assertNotEqual(5, methodToTest(n, m));}
  \choice \texttt{5 == methodToTest(n, m);}
  \choice \texttt{methodToTest(n, m).equals(5);}
\end{choices}
\begin{solution}
  Matches the source quiz key. In real JUnit~4 you typically call the method under test on an instance (\texttt{new Addition().methodToTest(n, m)}) or use a static import for \texttt{assertEquals}.
\end{solution}

\question[4]
\textit{(Select all that apply.)} Identify \emph{all} of the Java or OOP concepts being used here: \texttt{List mylist = new ArrayList<>();}
\begin{checkboxes}
  \CorrectChoice Generics
  \choice Encapsulation
  \CorrectChoice Constructor
  \choice Wrapper class
\end{checkboxes}
\begin{solution}
  The diamond \texttt{<>} and type arguments involve \textbf{generics}; \texttt{new ArrayList<>()} invokes a \textbf{constructor}. No new encapsulation boundary or wrapper conversion is shown by that line alone.
\end{solution}

\newpage

\question[4]
Identify the correct type of data structure to use in the code below:
\begin{lstlisting}[style=java]
____ countOfWords = getCounts();
Set keys = countOfWords.keySet();
for (String k : keys) {
    System.out.println(countOfWords.get(k));
}
\end{lstlisting}
\begin{choices}
  \choice \texttt{Set}
  \choice \texttt{List}
  \choice \texttt{Queue}
  \CorrectChoice \texttt{Map}
\end{choices}
\begin{solution}
  \texttt{keySet()} and \texttt{get(k)} are \texttt{Map} APIs (raw types here as in the source snippet; prefer \texttt{Map<String, Integer>} in real code).
\end{solution}

\question[4]
What would be the best collection to use in the code below if we want to ensure that no name appears twice when we print them out? (Assume \texttt{getNames()} returns the correct collection type.)
\begin{lstlisting}[style=java]
_____ names = getNames();
names.add("Bob");
names.add("Julie");
for (String name : names) {
    System.out.println(name);
}
\end{lstlisting}
\begin{choices}
  \CorrectChoice \texttt{Set}
  \choice \texttt{List}
  \choice \texttt{Queue}
  \choice \texttt{Map}
\end{choices}
\begin{solution}
  A \texttt{Set} enforces uniqueness (for a proper \texttt{equals}/\texttt{hashCode} contract on the elements). A \texttt{List} allows duplicates.
\end{solution}

\newpage

\question[4]
Identify the proper use of the interface declared here:
\begin{lstlisting}[style=java]
public interface Driveable {
    public void drive();
}
\end{lstlisting}
\begin{choices}
  \choice
  \begin{minipage}[t]{0.88\linewidth}
\begin{lstlisting}[style=java]
public class Car extends Driveable {
    public void drive() {
        System.out.println("driving a car...");
    }
}
\end{lstlisting}
  \end{minipage}
  \choice
  \begin{minipage}[t]{0.88\linewidth}
\begin{lstlisting}[style=java]
public class Car {
    @Inherit(Driveable.class)
    public void drive() {
        System.out.println("driving a car...");
    }
}
\end{lstlisting}
  \end{minipage}
  \CorrectChoice
  \begin{minipage}[t]{0.88\linewidth}
\begin{lstlisting}[style=java]
public class Car implements Driveable {
    public void drive() {
        System.out.println("driving the car...");
    }
}
\end{lstlisting}
  \end{minipage}
  \choice
  \begin{minipage}[t]{0.88\linewidth}
\begin{lstlisting}[style=java]
public class Car implements Driveable {
    public void stop() {
        System.out.println("stopping the car...");
    }
}
\end{lstlisting}
  \end{minipage}
\end{choices}
\begin{solution}
  Classes \textbf{implement} interfaces (\texttt{implements}) and must provide all abstract interface methods (here, \texttt{drive}). \texttt{extends} is for classes; \texttt{@Inherit} is not standard Java.
\end{solution}

\newpage

\question[4]
Can variable \texttt{y} be used at \texttt{B}? Why or why not? Can variable \texttt{x} be used at \texttt{A}? Why or why not?
\begin{lstlisting}[style=java]
public class Scopes {
    public void doStuff() {
        int x = 1;
        for (int i = 0; i < 5; i++) {
            int y = 1;
            // A
        }
        // B
    }
}
\end{lstlisting}
\begin{choices}
  \choice \texttt{x} cannot be used at \texttt{A}; \texttt{y} cannot be used at \texttt{B} because the variables are in different scopes
  \CorrectChoice \texttt{x} can be used at \texttt{A} because the block scope is within the method scope and variables from an enclosing scope can always be accessed; \texttt{y} cannot be used at \texttt{B} because variables in block scope cannot be accessed outside of the block
  \choice \texttt{x} cannot be used at \texttt{A} because the method scope is within the block scope; \texttt{y} can be used at \texttt{B} because variables in block scope can be accessed outside of the block
  \choice Both \texttt{x} and \texttt{y} can be used at both \texttt{A} and \texttt{B} because variables are always usable, regardless of scope
\end{choices}
\begin{solution}
  Inner blocks can see outer locals (\texttt{x} at \texttt{A}). \texttt{y} dies at the end of the loop body, so it is not in scope at \texttt{B}.
\end{solution}

\question[4]
Identify the correct output of this Java code and the correct explanation:
\begin{lstlisting}[style=java]
String s1 = "hello";
String s2 = " world";
s1.concat(s2);
System.out.println(s1);
String s3 = s1.concat(s2);
System.out.println(s3);
\end{lstlisting}
\begin{choices}
  \CorrectChoice \texttt{hello} then \texttt{hello world} --- strings are immutable, so \texttt{concat} does not modify \texttt{s1} but returns a new string
  \choice \texttt{hello world} twice --- strings are mutable, so \texttt{concat} appends to \texttt{s1} both times
  \choice \texttt{hello} twice --- strings cannot change and reference variables cannot be reassigned, so only \texttt{hello} prints twice
  \choice \texttt{hello world} then \texttt{hello world world} --- line~3 reassigns \texttt{s1}, and line~6 adds an extra \texttt{" world"}
\end{choices}
\begin{solution}
  \texttt{s1} stays \texttt{"hello"} until reassigned; only \texttt{s3} receives \texttt{"hello world"} from the second \texttt{concat}.
\end{solution}

\newpage

\question[4]
Identify the correct output of this code snippet:
\begin{lstlisting}[style=java]
public class Example {
    public static int countStatic = 0;
    public int count = 0;
    public static void main(String[] args) {
        Example e1 = new Example();
        Example e2 = new Example();
        e1.increment();
        e2.increment();
        System.out.println(e1.count);
        System.out.println(e2.count);
        System.out.println(countStatic);
    }
    public void increment() {
        count++;
        countStatic++;
    }
}
\end{lstlisting}
\begin{choices}
  \choice \texttt{2 2 2}
  \choice \texttt{1 1 1}
  \choice \texttt{2 2 1}
  \CorrectChoice \texttt{1 1 2}
\end{choices}
\begin{solution}
  Each instance has its own \texttt{count}; both increments run once per instance $\Rightarrow$ \texttt{1} and \texttt{1}. \texttt{countStatic} is shared $\Rightarrow$ \texttt{2}.
\end{solution}

\newpage

\question[4]
Identify the class that declares a constructor that allows for this object instantiation:
\begin{lstlisting}[style=java]
Example myExample = new Example("foo");
\end{lstlisting}
\begin{choices}
  \CorrectChoice
  \begin{minipage}[t]{0.88\linewidth}
\begin{lstlisting}[style=java]
public class Example {
    public String name;
    public Example(String s) {
        this.name = s;
    }
}
\end{lstlisting}
  \end{minipage}
  \choice
  \begin{minipage}[t]{0.88\linewidth}
\begin{lstlisting}[style=java]
public class Example {
    public String name;
    public Example() {
        System.out.println("foo");
    }
}
\end{lstlisting}
  \end{minipage}
  \choice
  \begin{minipage}[t]{0.88\linewidth}
\begin{lstlisting}[style=java]
public class Example {
    public String name;
    public void createExample(String s) {
        this.name = s;
    }
}
\end{lstlisting}
  \end{minipage}
  \choice
  \begin{minipage}[t]{0.88\linewidth}
\begin{lstlisting}[style=java]
public class Example {
    public String name;
    public static String createExample(String s) {
        return s;
    }
}
\end{lstlisting}
  \end{minipage}
\end{choices}
\begin{solution}
  \texttt{new Example("foo")} requires a constructor \texttt{Example(String)}.
\end{solution}

\newpage

\question[4]
Identify the valid use of method overloading.
\begin{choices}
  \CorrectChoice
  \begin{minipage}[t]{0.88\linewidth}
\begin{lstlisting}[style=java]
public class Example {
    public void doStuff() { System.out.println("hello"); }
    public void doStuff(int x) { System.out.println(x); }
}
\end{lstlisting}
  \end{minipage}
  \choice
  \begin{minipage}[t]{0.88\linewidth}
\begin{lstlisting}[style=java]
public class Example {
    public void doStuff() {}
    public String doStuff() {
        return "hello";
    }
}
\end{lstlisting}
  \end{minipage}
  \choice
  \begin{minipage}[t]{0.88\linewidth}
\begin{lstlisting}[style=java]
public class Example {
    public void doStuff() {
        System.out.println("hello");
    }
}
public class Sub extends Example {
    public void doOtherStuff(int x) {
        System.out.println(x);
    }
}
\end{lstlisting}
  \end{minipage}
  \choice
  \begin{minipage}[t]{0.88\linewidth}
\begin{lstlisting}[style=java]
public class Example {
    public void doStuff() {
        System.out.println("hello");
    }
}
public class Sub extends Example {
    public void doStuff() {
        System.out.println("goodbye");
    }
}
\end{lstlisting}
  \end{minipage}
\end{choices}
\begin{solution}
  Overloading: same method name, different parameter lists in one class. Duplicate signatures differing only by return type (second option) do not compile. The last option is overriding, not overloading.
\end{solution}

\newpage

\uplevel{\par\textbf{References, packages, and visibility (practice bank)}\par}

\question[4]
How many different objects and reference variables are created in the following code?
\begin{lstlisting}[style=java]
public void createObjects() {
    Object o1 = new Object();
    Object o2 = o1;
    Object o3 = new Object();
    o3 = o1;
}
\end{lstlisting}
\begin{choices}
  \choice 3 objects, 2 reference variables
  \CorrectChoice 2 objects, 3 reference variables
  \choice 3 objects, 4 reference variables
  \choice 1 object, 3 reference variables
\end{choices}
\begin{solution}
  Two \texttt{new Object()} calls $\Rightarrow$ two heap objects. \texttt{o2} and \texttt{o3} (after reassignment) alias \texttt{o1}'s object. Three local reference variables \texttt{o1}, \texttt{o2}, \texttt{o3}.
\end{solution}

\question[4]
Which access modifier should you use here to give access to the \texttt{String} field to class \texttt{B} but \emph{not} to class \texttt{C}?
\begin{lstlisting}[style=java]
// filename: A.java
package foo;
public class A {
    ____ String s = "hello";
}
// filename: B.java
package bar;
public class B extends A {
}
// filename: C.java
package bar;
public class C {}
\end{lstlisting}
\begin{choices}
  \choice \texttt{public}
  \choice \texttt{private}
  \CorrectChoice \texttt{protected}
  \choice Default (package-private; no access modifier)
\end{choices}
\begin{solution}
  \texttt{protected} allows the subclass \texttt{B} (different package) to access the field, while unrelated \texttt{C} in \texttt{bar} cannot. \texttt{public} would expose \texttt{C}; \texttt{private} would block \texttt{B}; default (package) would not allow \texttt{B} in \texttt{bar} to see \texttt{foo}'s field.
\end{solution}

\newpage

\uplevel{\par\textbf{Design patterns}\par}

\question[4]
What is the Factory design pattern?
\begin{solutionorbox}[1.25in]
  Encapsulates object creation behind a method or factory class so callers depend on interfaces/supertypes and creation rules stay in one place (variants: simple factory, factory method, abstract factory).
\end{solutionorbox}

\question[4]
What is the Singleton design pattern?
\begin{solutionorbox}[1.25in]
  Ensures a class has at most one instance and provides global access; useful for shared resources but easy to misuse (testing, threading)---often replaced with DI-managed beans.
\end{solutionorbox}

\endgradingrange{csoo}
